\subsection*{7. Starting Broad Can Be Rational in One Sense and Disastrous in Another}

Exploring broadly is often good. Operating broadly is often expensive. This distinction matters.

At the discovery stage, a founder might speak with learners, colleges, enterprise teams, and professional associations. That is not necessarily a lack of focus; it can be an efficient search for evidence. It becomes a problem when the founder builds 4 products, launches 4 marketing efforts, hires for 4 delivery models, and promises all 4 will be supported without having proven 1.

\textbf{Breadth in inquiry; concentration in commitment.} Before selecting a beachhead, use divergent research. After selecting one, concentrate the actual product and go-to-market resources while maintaining a limited ability to detect contrary evidence.

A company can also serve multiple customers because a single transaction structurally requires several sides: for example, a marketplace must coordinate supply and demand. In such cases, focused \emph{market density} matters more than choosing 1 side and ignoring the other.

\subsection*{8. The Beachhead Is a Hypothesis, Not a Marriage}

Founders sometimes resist narrowing because it feels like surrendering the rest of the market. The anxiety is understandable. A beachhead requires accepting that a specific group might not buy and that a rejected broader group might later become attractive.

A beachhead choice should be a \textbf{time-bounded allocation decision.} It states: ``For this learning period, we will pursue this customer situation intensely enough to learn whether the business can work.'' If evidence invalidates the choice, the founder changes it.

The cost of a narrow experiment is limited if fixed assets, long contracts, and reputation claims remain modest. The cost of building an organization for everyone before discovering demand can be ruinous. The rational founder should value reversibility, not avoid deciding.

\subsection*{9. Not Every Business Should Segment in the Same Dimension}

For a physical service, geography may be the correct axis: a dental clinic needs customers who can travel. For an enterprise security tool, regulation and workflow may matter more than city. For education, prior capability, desired transformation, schedule, funding source, and employer context may matter more than age. For a consumer snack, purchase occasion, retail channel, and sensory preference may dominate professional identity.

An effective beachhead uses the \textbf{dimension on which customer need, purchase behavior, and delivery requirements actually become similar.} Demographic similarity without those mechanisms is usually too weak.

\subsection*{10. Why ``Everyone Who Might Benefit'' Is an Impossible Target}

A product can be useful to a million people and purchasable by very few. Benefiting, wanting, paying, and successfully using are different events. If a founder says ``all developers can benefit from AI training,'' that does not specify which developers believe they need training now, which have access to money, which will devote weekly hours, which recognize the founder as credible, or which will accept the work required.

The beachhead restricts the proposition to a setting in which these conditions are more likely to align. It does not deny that outsiders might benefit. It identifies whom the business will organize itself to serve first.

\subsection*{11. A Useful Contrarian Test: Why Has the Incumbent Left This Opening?}

When you find an underserved group, do not assume you have discovered an irrational oversight. The incumbent may have made a rational choice. The segment might be hard to reach, require expensive support, pay little, produce high complaints, or conflict with the company's existing distribution model.

Ask whether the omission reflects \textbf{incumbent constraints} or \textbf{unattractive economics.} A software company may avoid a narrow compliance workflow because it needs expensive maintenance. A training company may avoid advanced technical feedback because expert reviewers cost more than mass recorded instruction. A college may not buy an external applied course because timetable and credit rules make integration difficult.

Your advantage may be that a smaller company can operate differently at a cost structure or specificity the incumbent cannot justify. But if you cannot make better economics or deliver a more valued result, neglected customers may simply remain unprofitable customers.

\subsection*{12. The Chapter's Governing Choice}

Selecting a beachhead involves 6 judgments:

\begin{enumerate}
    \item \textbf{Demand:} Does the group experience a sufficiently consequential, recurring, or time-sensitive job?
    \item \textbf{Coherence:} Can 1 proposition meaningfully serve multiple members of the group?
    \item \textbf{Access:} Can this founder repeatedly reach and earn consideration from them?
    \item \textbf{Advantage:} Is there a credible reason a small entrant could win here?
    \item \textbf{Economics:} Could sales and delivery produce acceptable contribution and cash flow?
    \item \textbf{Propagation:} Would success produce capabilities or proof that help the next group become accessible?
\end{enumerate}

These are not interchangeable. A beautiful long-term expansion path does not redeem a first market that cannot pay. Strong immediate margins do not excuse a product that requires impossible founder hours. Demand matters; access and delivery determine whether the demand can become this company's revenue.
